% Arte editorial de Molambudos. Compilar com XeLaTeX; parâmetros em variante.tex.
\def\Pages{436}
\def\Spine{25.99161680}
\def\CoverWidth{352.34161680}

\documentclass{article}
\usepackage[paperwidth=\CoverWidth mm,paperheight=236.35mm,margin=0pt]{geometry}
\usepackage{fontspec,graphicx,tikz,xcolor}
\setmainfont{EB Garamond}
\setsansfont{Lato}
\usepackage[pdfencoding=auto,unicode]{hyperref}
\hypersetup{pdftitle={Molambudos: capa completa, \Pages\ páginas},pdfauthor={Marcelo Dias de Carvalho Filho},pdfsubject={Brochura 160 x 230 mm; cor premium em papel branco; sangria 3,175 mm}}
\pagestyle{empty}
\definecolor{coal}{HTML}{010100}
\definecolor{backcoal}{HTML}{11120D}
\definecolor{gold}{HTML}{C3A574}
\definecolor{softgold}{HTML}{8E7955}
\definecolor{ivory}{HTML}{EDE1C9}
\begin{document}
\null
\begin{tikzpicture}[remember picture,overlay]
\begin{scope}[shift={(current page.north west)},x=1mm,y=-1mm]
\pgfmathsetmacro{\FrontX}{163.175+\Spine}
\pgfmathsetmacro{\FrontCenter}{\FrontX+80}
\pgfmathsetmacro{\SpineCenter}{163.175+\Spine/2}
\fill[backcoal] (0,0) rectangle (163.175,236.35);
\fill[coal] (163.175,0) rectangle (\CoverWidth,236.35);

% CONTRACAPA. A arte fica a 302,4 dpi efetivos, sem reamostragem.
\begin{scope}
\clip (0,0) rectangle (163.175,236.35);
\node[anchor=north west,inner sep=0] at (3.175,108) {\includegraphics[width=86mm]{arte_contracapa_original.png}};
\end{scope}
\draw[softgold,line width=.3pt] (19.175,19) -- (147.175,19);
\node[anchor=north west,inner sep=0,text=gold,align=left] at (19.175,28) {\fontsize{20}{22}\selectfont HÁ HISTÓRIAS QUE\vphantom{Á}\\\fontsize{20}{22}\selectfont O TEMPO ENTERRA.\\\fontsize{20}{22}\selectfont MAS NÃO APAGA.};
\node[anchor=north west,inner sep=0,text=ivory,text width=126mm,align=left,font=\fontsize{12.5}{16}\selectfont] at (19.175,63) {\hyphenpenalty=10000\exhyphenpenalty=10000\raggedright
Em 1915, a seca arranca Joaquim Antônio Correia de seu sertão. Anos depois, reduzido ao número 1.260 em uma instituição psiquiátrica, ele escreve para não desaparecer. Suas memórias atravessam a fome, o confinamento e os silêncios que nenhum prontuário consegue explicar.

\vspace{3mm}
Décadas mais tarde, a psiquiatra forense Lúcia Mendes reúne diários, laudos, cartas e fotografias para reconstituir uma história que parece resistir ao esquecimento. Mas cada fragmento abre outro caminho, e o arquivo começa a interpelar quem o lê.

\vspace{3mm}
Em \textit{Molambudos}, o horror nasce da memória ferida — e continua onde termina a página.
};
\node[anchor=north west,inner sep=0,text=gold,align=left,font=\fontsize{15}{18}\selectfont\itshape] at (19.175,158) {O ciclo recomeça\\a cada leitor.};
\node[anchor=north west,inner sep=0,text=ivory,align=left,font=\sffamily\fontsize{7.3}{10}\selectfont] at (19.175,180) {ROMANCE BRASILEIRO\\HORROR PSICOLÓGICO};
% Reserva silenciosa para código automático: 2 x 1,2 pol.; 6,35 mm do corte e da lombada.
\fill[white] (106.025,196.345) rectangle (156.825,226.825);

% LOMBADA. Toda a tipografia tem folga superior a 1,6 mm em relação às dobras.
\draw[softgold,line width=.35pt] (\SpineCenter,18) -- (\SpineCenter,30);
\node[rotate=-90,inner sep=0,text=gold] at (\SpineCenter,82) {\fontsize{22}{24}\selectfont\bfseries MOLAMBUDOS};
\node[rotate=-90,inner sep=0,text=ivory] at (\SpineCenter,158) {\fontsize{9.5}{12}\selectfont MARCELO DIAS DE CARVALHO FILHO};
\node[inner sep=0,text=gold,align=center,font=\sffamily\fontsize{7.3}{9}\selectfont] at (\SpineCenter,213) {OPENCODE\\BOOKS};

% FRENTE. Ilustração em prancha editorial, conservando resolução nativa >= 300 dpi.
\draw[softgold,line width=.3pt] (\FrontX+14,17) -- (\FrontX+146,17);
\node[inner sep=0,text=gold] at (\FrontCenter,33) {\resizebox{136mm}{!}{\bfseries MOLAMBUDOS}};
\node[inner sep=0,text=ivory] at (\FrontCenter,50) {\fontsize{16}{19}\selectfont O Diário do Paciente 1.260};
\node[inner sep=0,text=softgold] at (\FrontCenter,62) {\sffamily\addfontfeatures{LetterSpace=14}\fontsize{7.8}{10}\selectfont UM ROMANCE EM FRAGMENTOS};
\begin{scope}
\clip (\FrontCenter-43,73) rectangle (\FrontCenter+43,169);
\node[anchor=north,inner sep=0] at (\FrontCenter,40) {\includegraphics[width=86mm]{arte_frontal_original.png}};
\end{scope}
\draw[softgold,line width=.35pt] (\FrontCenter-44,72) rectangle (\FrontCenter+44,170);
\draw[gold,line width=.4pt] (\FrontCenter-46,72) -- (\FrontCenter-46,83);
\draw[gold,line width=.4pt] (\FrontCenter+46,159) -- (\FrontCenter+46,170);
\node[inner sep=0,text=ivory,align=center,font=\fontsize{12.5}{15}\selectfont\itshape] at (\FrontCenter,184) {Ela não é um monstro.\\Ela é uma consequência.};
\node[inner sep=0,text=gold] at (\FrontCenter,202) {\fontsize{13.4}{16}\selectfont Marcelo Dias de Carvalho Filho};
\node[inner sep=0,text=ivory] at (\FrontCenter,211) {\fontsize{10}{13}\selectfont\itshape Coautor e ilustrador: Marcelo Claro};
\draw[softgold,line width=.3pt] (\FrontX+14,220) -- (\FrontX+146,220);
\node[inner sep=0,text=softgold] at (\FrontCenter,225) {\sffamily\addfontfeatures{LetterSpace=12}\fontsize{7.2}{9}\selectfont OPENCODE BOOKS};
\end{scope}
\end{tikzpicture}
\end{document}
